\section{Free states and collision trees}
\label{sec:free-state-trees}\label{sec:free-states}

\subsection{The state left by a collision}

Suppose two insertions have a contraction \((z_1-z_2)^{-1}\).
After merging the first insertion with a nearby point, its contraction
with a third insertion is
\begin{equation}\label{fs:eq:first-expansion}
 \frac1{x+t}=\frac1x-\frac{t}{x^2}+\frac{t^2}{x^3}-\cdots.
\end{equation}
Keeping only the state at \(t=0\) removes every later coefficient.
The state left by the first collision must contain jets recording
these powers of \(t\).  There is a second issue.  An order parameter
can carry a degree-one differential form, so the contraction and the
state must be composed with the tensor signs of complexes.

We give the construction for finite polynomial inputs.  The ground
field \(k\) is arbitrary.  Complexes are cohomological, and a degree
\(d\) map raises degree by \(d\).  For homogeneous maps and vectors,
\[
 (f\otimes g)(u\otimes v)=(-1)^{|g||u|}f(u)\otimes g(v).
\]
All signs mean their images in \(k\).  In characteristic two we
still impose that the square of every exterior generator is zero.
Coordinates and scale variables have degree zero and differential zero.

The result concerns a free state algebra with a specified contraction.
It does not begin with a bulk action or a boundary condition.
A physical realization would have to supply the stated states,
propagator, and differential.  The distinction will matter when the
coefficient complex changes characteristic.

\subsection{Jets, translations, and the coefficient algebra}

Fix a finite set of species \(a=1,\ldots,d\).  The state algebra is
\begin{equation}\label{fs:eq:state-algebra}
 S=k[x_{a,r}:r\ge0]\otimes\Lambda(\eta^a_r:r\ge0),
 \qquad |x_{a,r}|=0,\quad |\eta^a_r|=1,\quad d_S=0.
\end{equation}
Its elements are finite polynomials.  The exterior algebra imposes
\(\eta\eta'=-\eta'\eta\) and \(\eta^2=0\).
The auxiliary weight of both \(x_{a,r}\) and \(\eta^a_r\) is
\(r+1\).  Thus \(S=\bigoplus_{E\ge0}S_E\), with each \(S_E\)
finite-dimensional.  Indeed, weight \(E\) uses at most \(E\)
factors and only jets of index less than \(E\).
This auxiliary grading differs from a conformal grading in which
\(x_{a,0}\) has weight zero.  In the latter grading the independent
states \(1,x_{a,0},x_{a,0}^2,\ldots\) all have weight zero, so
finite-dimensionality of homogeneous pieces would fail.

For a formal variable \(v\), define an algebra map
\begin{equation}\label{fs:eq:hasse-translation}
 \tau_v(x_{a,r})=\sum_{m\ge0}\binom{r+m}{r}x_{a,r+m}v^m,
 \qquad
 \tau_v(\eta^a_r)=\sum_{m\ge0}\binom{r+m}{r}\eta^a_{r+m}v^m.
\end{equation}
The sums are interpreted weight by weight.  Every coefficient is
finite, and the exterior relations are preserved.  No division in
\(k\) is used.

\begin{lemma}[Translation and the jet kernel]\label{fs:lem:hasse}
The maps satisfy \(\tau_v\tau_w=\tau_{v+w}\) and \(\tau_0=\id\).
For \(r,s\ge0\), put
\[
 K_{rs}(u)=(-1)^r\binom{r+s}{r}u^{-r-s-1}.
\]
In \(k[u,u^{-1}][[v,w]]\),
\begin{equation}\label{fs:eq:kernel-translation}
 K_{rs}(u+v-w)=
 \sum_{m,n\ge0}\binom{r+m}{r}\binom{s+n}{s}
 K_{r+m,s+n}(u)v^mw^n.
\end{equation}
Also \(K_{sr}(-u)=-K_{rs}(u)\).
\end{lemma}

\begin{proof}
The coefficient of \(v^mw^n x_{a,r+m+n}\) in the composite
translation is
\(\binom{r+n}{r}\binom{r+n+m}{r+n}\).
The coefficient in \(\tau_{v+w}\) is
\(\binom{r+m+n}{r}\binom{m+n}{m}\).
Both integers count an ordered division of \(r+m+n\) objects
into parts of sizes \(r,m,n\).  The same argument applies to
\(\eta\), and multiplicativity proves the translation law.

Define a Hasse derivative by
\(\partial_z^{[r]}F(z)=[v^r]F(z+v)\).
On a monomial it is \(\binom nr z^{n-r}\), with the integer
binomial convention also for negative \(n\).  Coefficient extraction
gives
\(\partial^{[r]}\partial^{[m]}=\binom{r+m}{r}\partial^{[r+m]}\).
The geometric expansion of \((u+v-w)^{-1}\) has coefficient
\((-1)^r\binom{r+s}{r}u^{-r-s-1}\) at \(v^rw^s\).
Apply the two Hasse derivatives to that expansion and use their
composition law.  This proves \eqref{fs:eq:kernel-translation} over
the integers, hence over \(k\).  Substitution of \(-u\) proves
the last identity.
\end{proof}

Let \(I=\{1,\ldots,n\}\).  Write
\[
 R_I=k[z_i:i\in I],\qquad
 L_I=R_I[(z_i-z_j)^{-1}:i<j].
\]
The polynomial differential forms on a simplex with vertex set \(V\)
are the algebra
\[
 \Omega(V)=
 k[t_v:v\in V]\otimes\Lambda(dt_v:v\in V)
 /(\sum t_v-1,\sum dt_v),
 \qquad d(t_v)=dt_v.
\]
The differential follows the graded product rule.  A face restriction
sets the missing \(t_v,dt_v\) to zero.

Take \(V=\Ord(I)\), the set of all total orders on \(I\).
For a nonempty face \(J\subseteq\Ord(I)\), let \(R_J\) invert
precisely the differences whose comparisons vary among orders in \(J\).
Define \(C_I\subseteq L_I\otimes\Omega(\Ord(I))\) by requiring
that restriction to each face have coefficients in \(R_J\).
This is a differential graded algebra over any field.  We will test
separately whether it computes gluing cohomology.

The degree-zero order function and its kernel are
\begin{equation}\label{fs:eq:order-kernel}
 q_{ij}=\sum_{\sigma\in\Ord(I)}t_\sigma\,1_{i<_\sigma j},
 \qquad c_{ij}^{rs}=dq_{ij}\,K_{rs}(z_i-z_j).
\end{equation}
On a face fixing the comparison, \(dq_{ij}=0\), so this formula
means zero without inverting that difference.  On every other face
the inverse is allowed.  Thus \(c_{ij}^{rs}\in C_I^1\), it is
closed, and \(c_{ji}^{sr}=c_{ij}^{rs}\).
For \(J\subseteq I\), forgetting labels gives a coefficient map
\(C_J\to C_I\): it retains the coordinates labelled by \(J\)
and sends a barycentric coordinate \(t_\tau\) to
\(\sum_{\sigma|_J=\tau}t_\sigma\).
This preserves the sum-one relation and every face condition, because
a comparison in \(J\) varies exactly when the same comparison varies
among the restricted orders.  Relabellings act by the corresponding
coordinate and barycentric permutations.

\subsection{Finite contractions as chain maps}

The left derivative \(\partial_{\eta}\) deletes an occurrence of
\(\eta\) with sign \((-1)^b\), where \(b\) counts the odd
factors preceding it.  It has degree \(-1\), squares to zero,
and satisfies the graded product rule.  Derivatives in even variables
are ordinary polynomial derivatives of degree zero.

On \(S^{\otimes I}\), put
\begin{equation}\label{fs:eq:pair-operators}
 P_{ij}^{rs}=\sum_a\left(
 \partial_{x_{a,r}^{(i)}}\partial_{\eta_s^{a(j)}}+
 \partial_{\eta_r^{a(i)}}\partial_{x_{a,s}^{(j)}}\right),
 \qquad
 D_{ij}=\sum_{r,s\ge0}c_{ij}^{rs}\otimes P_{ij}^{rs}.
\end{equation}
It is useful to separate the odd rational operator from its one-form:
\begin{equation}\label{fs:eq:rational-pair-operator}
 \mathsf D_{ij}=\sum_{r,s\ge0}K_{rs}(z_i-z_j)P_{ij}^{rs},
 \qquad D_{ij}=dq_{ij}\otimes\mathsf D_{ij}.
\end{equation}
Composition acts from right to left.  Both \(P_{ij}^{rs}\) and
\(\mathsf D_{ij}\) have degree \(-1\), while \(D_{ij}\) has
degree zero.  On a fixed polynomial input all
these sums are finite.  On coefficient-valued inputs the convention is
\[
 (\alpha\otimes P)(\beta\otimes v)
 =(-1)^{|P||\beta|}\alpha\beta\otimes P(v).
\]

\begin{lemma}[Commuting square-zero pair operations]\label{fs:lem:square-zero}
One has \(D_{ij}^2=0\), \(D_{ij}D_{kl}=D_{kl}D_{ij}\), and
\([d,D_{ij}]=0\).  These identities hold in every characteristic.
\end{lemma}

\begin{proof}
Each summand of a \(P\) contains one odd derivative and commuting
even derivatives.  Two such summands anticommute.  If they use the
same odd derivative their product is zero.  Hence all the \(P\)'s
anticommute, including shared-label cases.
All coefficients of \(D_{ij}\) contain the same one-form
\(dq_{ij}\), whose square is zero.  This proves \(D_{ij}^2=0\).
For different pairs, both the coefficient one-forms and the odd
operators anticommute.  The tensor composition contributes one
additional minus sign in either order.  The two products are equal.
The coefficients are closed and the state differential is zero,
which proves the chain identity.  The argument uses exterior
relations, so it also applies in characteristic two.
\end{proof}

Let \(E_z:S^{\otimes I}\to\prod_E(R_I\otimes S_E)\) translate
the \(i\)-th factor by \(\tau_{z_i}\), then multiply in label order.
The Koszul sign appears whenever odd factors are reordered.
For a pole tuple \(N=(N_{ij})\), set
\[
 D_I(N)=\left(\prod_{i<j}(z_i-z_j)^{-N_{ij}}\right)R_I.
\]
The target \(\mathscr M_I\) consists of elements of
\(\prod_E(C_I\otimes S_E)\) whose coefficients lie in
\(D_I(N)\otimes\Omega(\Ord(I))\) for one common finite tuple \(N\).
The product is over the auxiliary state weight, with coefficient
spaces discrete.  Thus state weights may grow, but a single input
has a finite common denominator before collision expansion.

\begin{theorem}[The free state operation]\label{fs:thm:free-operation}
The formula
\begin{equation}\label{fs:eq:free-operation}
 Y_I=E_z\prod_{i<j}(1+D_{ij}):S^{\otimes I}\longrightarrow\mathscr M_I
\end{equation}
defines a degree-zero chain map over every field.
The product is independent of its pair ordering.
It commutes with permutations with their Koszul signs, and deleting
a unit input deletes the corresponding label.  The unary operation
is \(\tau_z\), and the empty operation is the state \(1\).

For polynomial inputs, set \(N_{ij}\) equal to the largest
\(r+s+1\) among opposite-parity, same-species factors available
at sites \(i,j\), or zero if there is none.  Every output
coefficient has pole order at most \(N_{ij}\) on that diagonal.
\end{theorem}

\begin{proof}
The product contains finitely many pair factors.  Derivatives remove
factors, so every evaluation contains only finitely many contraction
choices.  A nonzero choice uses at most one contraction on each
pair.  Its denominator on that pair has exponent \(r+s+1\).
Translations of the remaining factors have polynomial coefficients
in each fixed state weight and introduce no poles.  This proves
the target and the stated bounds.
Lemma~\ref{fs:lem:square-zero} proves order independence and the chain
property.  Face restriction commutes with the construction: a forbidden
pair coefficient restricts to zero by \eqref{fs:eq:order-kernel}.
The remaining contractions have legal denominators on that face.

Relabelling sends the pair kernel to the same kernel and the odd
derivatives to their signed tensor permutations.  Multiplication in
the graded commutative state algebra has exactly these signs.
This proves permutation compatibility.  Every derivative of a unit
input vanishes, and \(\tau_z(1)=1\), proving the unit statements.
\end{proof}

In characteristic zero, the product is the exponential of
\(\sum D_{ij}\).  This follows by expanding the exponential and
using commuting square-zero factors.  Formula~\eqref{fs:eq:free-operation}
is the definition in arbitrary characteristic.

\begin{example}[A nonzero two-form and a vanishing jet kernel]
With two species, use the three inputs
\(x_{1,0},x_{2,0},\eta^1_0\eta^2_0\).
The scalar-state component is \(c_{13}^{00}c_{23}^{00}\).
Deleting \(\eta^2_0\) after the first contraction crosses its
coefficient one-form; moving the second one-form back to pair order
gives a second minus sign.  The resulting coefficient is positive.
On the face with orders \(123,312,132\), write the first two
barycentric coordinates as \(a,b\).  Then
\(q_{13}=1-b\), \(q_{23}=a\), and this component is
\[
 \frac{da\,db}{(z_1-z_3)(z_2-z_3)}.
\]
It is nonzero over every field and has exactly the two displayed
simple divisorial poles.  On the binary tree \((12)3\), both order
differentials become the outer \(dq\), and this component vanishes.

Individual jet contractions can instead vanish for arithmetic reasons.
For the binary inputs \(x_r,\eta_s\), the contracted component is
\(dq\,K_{rs}(u)\).  Its pole order is exactly \(r+s+1\)
if \(\binom{r+s}{r}\ne0\) in \(k\), and it is zero otherwise.
For instance \(K_{11}(u)=-2u^{-3}\) vanishes in characteristic two,
whereas \(K_{10}(u)=-u^{-2}\) survives in every characteristic.
\end{example}

\subsection{A common tree domain and its insertion identity}

A finite rooted tree has labelled leaves.  Each internal vertex has
a finite set of children, local child coordinates, and an order simplex
on those children.  Fix one child coordinate equal to zero at each
vertex.  For every nonroot internal vertex \(v\), introduce a
scale \(\epsilon_v\).  The physical displacement inside that vertex
is multiplied by the product of scales along the path from the root.
The coordinate of a leaf is the sum of its displacements along that
path.  For example,
\[
 z_{a,b,c}=Z_a+\epsilon_a U_{a,b}
                  +\epsilon_a\epsilon_{a,b}V_{a,b,c}.
\]
Leaf orders are formed by replacing each child in its local order
with the recursively ordered leaves of that subtree.

For chosen nonempty faces at all vertices, let \(B_T\) be the
tensor product over \(k\) of their coordinate overlap rings and
their polynomial-form algebras.  The common scale ring is
\begin{equation}\label{fs:eq:tree-ring}
 B_T[[\epsilon_v]][(\prod_v\epsilon_v)^{-1}].
\end{equation}
We then take a product over the state weights, retaining a common
finite lower bound in each scale for a fixed original polynomial input.
Coefficients at a specified scale multiindex and state weight are finite.

\begin{lemma}[Coordinate and order pullbacks]\label{fs:lem:tree-pullback}
The tree coordinate substitution sends every permitted rational
coefficient to \eqref{fs:eq:tree-ring}.
The leaf barycentric coordinate \(t_\sigma\) pulls back to the
product of the vertex barycentric coordinates producing \(\sigma\),
or to zero if no such choices exist.  Differentiation defines the
pullback on forms.  Under this map \(q_{ij}\) becomes the order
function of the two children at the lowest common ancestor of
leaves \(i,j\).
\end{lemma}

\begin{proof}
At their lowest common ancestor \(v\), the leaf difference has form
\[
 z_i-z_j=\lambda_v(u_a-u_b+h),
\]
where \(\lambda_v\) is the ancestor-scale monomial and \(h\)
has positive degree in the descendant scales.  If the leaf comparison
varies, the corresponding child comparison varies at \(v\), so
\(u_a-u_b\) is a unit in the chosen overlap ring.  Geometric
inversion gives a power series after the finite factor
\(\lambda_v^{-1}\).  Finitely many original denominators give
common finite lower bounds in all scales.

The vertex choices producing a compatible leaf order are unique,
by restriction to consecutive subtrees.  Summing their products
over all leaf orders gives the product of the vertex sums, which
is one.  Hence the barycentric relations and their differentials
are preserved.  In the sum defining \(q_{ij}\), the comparison
depends only on the two children at \(v\).  The sums over all
other choices are one, proving the last assertion.
\end{proof}

\begin{theorem}[Finite-tree state insertion]\label{fs:thm:tree-insertion}
Take polynomial leaf states in \(S\).  At each vertex apply
\eqref{fs:eq:free-operation} to its child states using their physical
relative displacements and that vertex's order functions.  Expand
all rational factors in the common ring \eqref{fs:eq:tree-ring}.
The resulting state-valued series equals the tree pullback of
\(Y_I\).  This equality holds in every characteristic and is
compatible with grafting finite trees.
The domain consists of series generated in this way from the fixed
polynomial leaf states, not arbitrary completed intermediate states.
\end{theorem}

\begin{proof}
Expand the finite pair product before applying translations.
A term chooses contracted pairs of individual factors and leaves
the other factors uncontracted.  Each factor can be used at most
once, because it is deleted by its derivative.  Multiplicities from
repeated even factors count the choices of their occurrences.
The derivative definition assigns the Koszul sign to this matching.

Partition a matching by the lowest common ancestor of each contracted
pair.  This assigns it uniquely to the vertices where the iterated
operation must perform those contractions.  Conversely, contractions
at vertices choose disjoint surviving factors and reconstruct the
leaf matching.  If two contractions use the same child pair at a
vertex, their common \(dq\) squares to zero.  The direct pullback
of that matching is also zero.  Thus the same terms vanish on
both sides.

For a contraction between translated child states, expand their
translations by \eqref{fs:eq:hasse-translation}.  Identity
\eqref{fs:eq:kernel-translation} sums exactly these jet contractions
to the expansion of the original leaf kernel.  It includes the
physical scale factors because the arguments are the actual relative
displacements.  The uncontracted factors agree by
\(\tau_v\tau_w=\tau_{v+w}\).
Every intermediate operation has total degree zero.  Moving its
coefficient form to the left contributes the tensor sign of its
odd deleted factors.  Hence the iterated sign is the sign obtained
by deleting the same matching in the original ordered tensor.
Equivalently, one may interchange successive pair contractions using
Lemma~\ref{fs:lem:square-zero}, which preserves their signs.

At a fixed state weight, finitely many uncontracted jet indices occur.
For contracted jets at a vertex with scale monomial \(\lambda\),
write a child displacement as \(\lambda v\).  The factor
\(\lambda^m v^m\) from its extra jet index cancels the factor
\(\lambda^{-m}\) from the outer jet kernel.  The remaining
\(v^m\) has positive descendant-scale degree.  A specified scale
coefficient therefore uses finitely many extra indices after the
fixed negative shifts from the original poles.  The denominator calculation in
Lemma~\ref{fs:lem:tree-pullback} gives those common lower bounds.
Thus the matching equality holds coefficientwise in the stated
common ring.  Grafting groups the same matchings by a longer path
of lowest common ancestors.  The same partition and translation
identities prove compatibility.
\end{proof}

For a single collision \(z_1=z_2+t\), put
\(A=N_{12}\), \(B_j=N_{1j}\), \(C_j=N_{2j}\).
Each state-weight coefficient has a polynomial numerator after its
common denominator is cleared.  Its coefficient of \(t^q\) is
zero for \(q<-A\).  For \(m=q+A\ge0\), it is a sum of terms
with merged pole exponents
\begin{equation}\label{fs:eq:state-pole-budget}
 B_j+C_j+h_j,\qquad \sum_jh_j\le m,\qquad h_j=0\text{ if }B_j=0.
\end{equation}
To prove this, expand the polynomial numerator and
\((x_j+t)^{-B_j}\).  A term of numerator degree \(\ell\)
and cross Taylor degrees \(h_j\) satisfies
\(\ell+\sum h_j=m\).  The coefficients are the integers
\((-1)^{h_j}\binom{B_j+h_j-1}{h_j}\), reduced in \(k\).
This proves the bound without a characteristic restriction.
Successive use gives a finite bound at every vertex of a fixed tree.
Cancellation among states or exterior forms can make it strict.

\subsection{Two three-input mergers and a chain homotopy}

Work on the fully separated rational carrier
\[
 B=k[t,x,t^{-1},x^{-1},(x+t)^{-1}]
       \otimes k[p,q]\otimes\Lambda(dp,dq).
\]
Its differential differentiates only \(p,q\).
Put \(z_1=x+t\), \(z_2=x\), \(z_3=0\), and write
\(X(v)=\tau_v(x_0)\), \(H(v)=\tau_v(\eta_0)\) for one species.
Use the product over state weights as above.
The order pullbacks for the two binary trees are
\[
 \begin{array}{c|ccc}
 &q_{12}&q_{13}&q_{23}\\\hline
 (12)3&p&q&q\\
 1(23)&q&q&p.
 \end{array}
\]
In both cases \(p\) orders the inner pair and \(q\) orders
the two outer blocks.  These are different maps of order parameters,
even before one chooses Laurent expansions.
The symbols \(Y_L,Y_R\) below denote rational representatives:
their expansions give the iterated operations in the two respective
tree rings.  Their comparison is made in \(B\), without identifying
the two Laurent-series domains.

\begin{example}[Three states with one odd factor]\label{fs:ex:three-even-even-odd}
For \(x_0\otimes x_0\otimes\eta_0\), the two operations are
\[
 \begin{aligned}
 Y_L&=X(x+t)X(x)\eta_0+
       \frac{dq}{x+t}X(x)+\frac{dq}{x}X(x+t),\\
 Y_R&=X(x+t)X(x)\eta_0+
       \frac{dq}{x+t}X(x)+\frac{dp}{x}X(x+t).
 \end{aligned}
\]
Thus they are generally unequal, but
\begin{equation}\label{fs:eq:first-homotopy}
 Y_R-Y_L=d\left(\frac{p-q}{x}X(x+t)\right).
\end{equation}
Both the pole and the primitive have exact order one along \(x=0\),
detected by their \(x_0\)-coefficient.  The primitive is a degree
\(-1\) operation on these degree-one inputs.
All formulas hold in characteristic two as well.
\end{example}

\begin{example}[The sign from two odd factors]\label{fs:ex:three-odd}
For \(\eta_0\otimes x_0\otimes\eta_0\), one obtains
\[
 \begin{aligned}
 Y_L&=H(x+t)X(x)\eta_0+\frac{dp}{t}\eta_0-
                                      \frac{dq}{x}H(x+t),\\
 Y_R&=H(x+t)X(x)\eta_0+\frac{dq}{t}\eta_0-
                                      \frac{dp}{x}H(x+t).
 \end{aligned}
\]
Deleting the last odd factor crosses the first odd factor and produces
the minus sign.  Their difference is
\[
 Y_R-Y_L=d\left((q-p)
                  \left(\frac{\eta_0}{t}+\frac{H(x+t)}{x}\right)\right).
\]
The two poles are simple and exact on their respective generic divisors.
In characteristic two the displayed minuses become pluses, while
the exterior squares remain zero.
\end{example}

\begin{proposition}[Jets cannot be discarded]\label{fs:prop:jet-obstruction}
Replacing \(\tau_v\) by the identity on \(k[x_0]\otimes\Lambda\eta_0\),
while retaining the kernel \(dq/(z_i-z_j)\), fails binary-tree insertion.
The failure occurs on the three nonunit inputs of
Example~\ref{fs:ex:three-even-even-odd} over every field.
\end{proposition}

\begin{proof}
Merge the first two inputs.  With no jets the inner output is
\(x_0^2\), and the outer binary formula is
\[
 x_0^2\eta_0+\frac{2dq}{x}x_0.
\]
The direct three-input formula pulled to the same order square is
\[
 x_0^2\eta_0+dq\left(\frac1{x+t}+\frac1x\right)x_0.
\]
After expansion at \(t=0\), its coefficient of \(t\) is
\(-dq\,x_0/x^2\), whereas the iterated expression has coefficient
zero.  This is nonzero over every field, including characteristic two.
In \(S\), the coefficient \(x_1\) of \(\tau_t(x_0)\) gives
the missing contraction \(K_{10}(x)=-x^{-2}\).
\end{proof}

\subsection{Integral higher comparisons on the separated carrier}

Let \(B_0\) be a commutative differential graded \(k\)-algebra
of order parameters and put \(B=L_I\otimes_k B_0\), with
differential zero on \(L_I\).  Choose two degree-zero order functions
\(q_e^0,q_e^1\in B_0\) for each pair \(e\); in particular,
they introduce no spatial poles.  No geometric meaning
of these arbitrary choices is assumed.  Set
\[
 L_e=\sum_{r,s}K_{rs}(z_i-z_j)P_{ij}^{rs},\qquad
 F_e^a=1+dq_e^a\otimes L_e,\qquad
 H_e=(q_e^1-q_e^0)\otimes L_e.
\]
Here \(L_e,H_e\) have degree \(-1\), and \(F_e^a\) has degree zero.
For an operator \(T\) of degree \(h\), write
\(\partial T=dT-(-1)^hTd\).  The product rule gives
\(\partial H_e=F_e^1-F_e^0\) and \(\partial F_e^a=0\).
These formulas are defined on polynomial state inputs.

\begin{theorem}[Cubical chain comparisons]\label{fs:thm:cubical}
Order the finitely many pairs as \(e_1,\ldots,e_M\).
The two operations \(E_z\prod_jF_{e_j}^0\) and
\(E_z\prod_jF_{e_j}^1\) are chain homotopic by
\begin{equation}\label{fs:eq:telescoping-homotopy}
 H=E_z\sum_{j=1}^M
 \left(\prod_{i<j}F_{e_i}^1\right)H_{e_j}
 \left(\prod_{i>j}F_{e_i}^0\right).
\end{equation}
All products and homotopies satisfy the same pairwise pole bounds
as the corresponding finite contraction terms.
There are compatible higher homotopies for every face of the
\(M\)-dimensional cube, over every field.
\end{theorem}

\begin{proof}
Apply \(\partial\) to \eqref{fs:eq:telescoping-homotopy}.
All \(F\)'s are closed and have degree zero.  The resulting sum
telescopes to
\(E_z(\prod F^1-\prod F^0)\), proving the chain-homotopy identity.
Each \(H_e\) has the same spatial kernel as its pair contraction.
Changing \(dq_e\) to \(q_e^1-q_e^0\) introduces no additional
spatial denominator.

For completeness, a cube face specifies some coordinates as endpoints
zero or one and a set \(J\) of varying coordinates.  Assign to it
the ordered product with factor \(H_{e_j}\) for \(j\in J\)
and factor \(F_{e_j}^{a_j}\) at an endpoint, followed by \(E_z\).
This map has degree \(-|J|\).  If
\(J=\{j_1<\cdots<j_r\}\), the product rule gives its differential
as the sum
\[
 \sum_{\nu=1}^r(-1)^{\nu-1}
 \bigl(\text{map on the face }a_{j_\nu}=1
       -\text{map on the face }a_{j_\nu}=0\bigr).
\]
This is the oriented boundary formula for the cube.
Its boundary squares to zero because the two orders of fixing a
pair of coordinates occur with opposite signs.
Hence the assigned maps are coherent higher chain homotopies.
No integration or division is used.
\end{proof}

The separated coefficient ring is essential here.  For example,
the primitive in \eqref{fs:eq:first-homotopy} has coefficient
\(1/x\) at the order-square vertex \((p,q)=(1,0)\).
At an order vertex the permitted coefficient ring contains only
regular coordinate functions.  Thus this homotopy need not preserve
the order-face restrictions.  A comparison after full localization
is an actual chain comparison, but it does not prove one before
localization.

\subsection*{Changing the coefficient complex}

The pair operation uses polynomial forms on order simplices.
Its formulas make sense over every field, but this does not identify
that coefficient complex with derived gluing coefficients.
The binary interval already distinguishes them: in characteristic
\(p>0\), the form \(u^{-1}ds\) is the differential of
\(u^{-1}(s-s^p)\), whose endpoint values are zero.
The overlap interval of Section~\ref{sb:sec:two-chart}
retains the nonzero class \([u^{-1}\theta]\), where \(\theta\)
is its degree-one overlap generator.
Proposition~\ref{sb:prop:coefficient-comparison} gives the precise comparison
and its characteristic-zero hypothesis.

\subsection{The remaining extension problem}

The finite-polynomial domain is part of the construction.
For the energy completion \(\widehat S=\prod_ES_E\), the states
\(\eta_r\) tend to zero as \(r\to\infty\).
Nevertheless, the scalar-form component of the binary operation on
\(x_0\otimes\eta_r\) is
\[
 \frac{dq}{u^{r+1}},
\]
because \(K_{0r}(u)=u^{-r-1}\).
These nonzero terms lie in the same discrete state-weight-zero
coefficient space.  They do not tend to zero.
Thus the operation has no continuous extension to all of
\(\widehat S\) with this topology.  A simultaneous pole and
state completion requires an additional continuity theorem.

For a general finite order atlas, one possible integral coefficient
carrier is the ordered overlap complex: degree \(r\) is the
direct sum of the coordinate rings on intersections of \(r+1\)
distinct charts, with alternating restriction differential.
Fixing an order on charts defines its cup product by multiplying
the first \(a+1\) and last \(b+1\) restrictions on an
\((a+b+1)\)-fold intersection.  The common chart occurs once.
Splitting the alternating omission sum at the common chart proves
the differential product rule: the two terms at that chart cancel
with opposite signs.  Cancellation of double omissions proves \(d^2=0\).
This associative complex is defined over every field,
but its cup product is not strictly graded commutative.

\begin{problem}[State operations on integral gluing coefficients]
Construct multilinear operations on the overlap complexes that retain
the binary class \([u^{-1}\theta]\), commute with Hasse
translation, and satisfy the finite-tree identities with their full
chain homotopies.  The maps must preserve every allowed overlap ring.
Any extension to completed state spaces must specify a joint
pole-state topology and prove continuity of every insertion.
\end{problem}

The free polynomial-form operations and the separated cubical
comparisons above are explicit chain-level constructions.
The lost-class calculation identifies why they cannot be transported
to the integral gluing carrier by a pole-preserving chain map from
the same polynomial-form model.  A construction on that corrected
carrier must retain the overlap class from the start.
